\lyseisE{Επανάληψη κεφαλαίου Γ3}
\solr{Γ3}{1} Πειράματα τύχης είναι τα α) και γ). Το β) έχει πάντα αποτέλεσμα $12$ και το δ) το διαβάζουμε από το ημερολόγιο.
\solr{Γ3}{2} Μήνες με $30$ ημέρες: $\{\text{Απρίλιος},\text{Ιούνιος},\text{Σεπτέμβριος},\text{Νοέμβριος}\}$.
\solr{Γ3}{3} Από τη ρίζα ξεκινούν $2$ κλαδιά (Κ, Γ) και από το καθένα $6$ (ζάρι). Αποτελέσματα $2\cdot6=12$. Με Κεφάλι και άρτιο: $(\mathrm{Κ},2),(\mathrm{Κ},4),(\mathrm{Κ},6)$, $3$ αποτελέσματα. Πιθανότητα $\dfrac{3}{12}=\dfrac{1}{4}=0{,}25=25\%$.
\solr{Γ3}{4} $3\cdot4\cdot2=24$ γεύματα.
\solr{Γ3}{5} Ισοπίθανα είναι τα $36$ ζευγάρια. α) Άθροισμα $10$: $(4,6),(5,5),(6,4)$, $\dfrac{3}{36}=\dfrac{1}{12}\approx8{,}3\%$. β) Άθροισμα $10$, $11$ ή $12$: $3+2+1=6$ ζευγάρια, $\dfrac{6}{36}=\dfrac{1}{6}\approx16{,}7\%$.
\solr{Γ3}{6} Σύνολο $10$ μπάλες. α) $\dfrac{3}{10}=0{,}3=30\%$. β) Όχι κόκκινη: $3+2=5$, $\dfrac{5}{10}=\dfrac{1}{2}$. γ) Δεν υπάρχει κίτρινη: $0$.
\solr{Γ3}{7} Κερδίζει: $\dfrac{5}{50}=\dfrac{1}{10}=0{,}1=10\%$. Δεν κερδίζει: $\dfrac{45}{50}=\dfrac{9}{10}=0{,}9=90\%$.
\solr{Γ3}{8} Σπαθί: $\dfrac{13}{52}=\dfrac{1}{4}$. Εξάρι: $\dfrac{4}{52}=\dfrac{1}{13}$.
\solr{Γ3}{9} α) Πολλαπλάσια του $3$: $3,6,9$, $\dfrac{3}{10}=0{,}3$. β) Πρώτοι: $2,3,5,7$, $\dfrac{4}{10}=\dfrac{2}{5}=0{,}4$. γ) Αδύνατο: $0$.
\solr{Γ3}{10} Φράουλα: $0{,}4\cdot30=12$. Λεμόνι: $\dfrac{1}{3}\cdot30=10$. Άλλες: $30-12-10=8$.
\solr{Γ3}{11} Οκτώ ισοπίθανα αποτελέσματα. Ακριβώς ένα Γράμματα: $\mathrm{ΓΚΚ},\mathrm{ΚΓΚ},\mathrm{ΚΚΓ}$, $\dfrac{3}{8}$. Και οι τρεις ίδιες: $\mathrm{ΚΚΚ},\mathrm{ΓΓΓ}$, $\dfrac{2}{8}=\dfrac{1}{4}$.
\solr{Γ3}{12} Σχετική συχνότητα $\dfrac{52}{300}\approx0{,}173$. Η πιθανότητα είναι $\dfrac{1}{6}\approx0{,}167$. Είναι πολύ κοντά.
\solr{Γ3}{13} Το αποτέλεσμα $\mathrm{ΓΚ}$ λείπει και τα τρία αποτελέσματα που έγραψε δεν είναι ισοπίθανα. Ο σωστός δειγματικός χώρος είναι $\{\mathrm{ΚΚ},\mathrm{ΚΓ},\mathrm{ΓΚ},\mathrm{ΓΓ}\}$. Ένα Κ και ένα Γ: $\mathrm{ΚΓ}$ και $\mathrm{ΓΚ}$, δηλαδή $\dfrac{2}{4}=\dfrac{1}{2}$.
\solr{Γ3}{14} Οι αριθμοί είναι $123,132,213,231,312,321$ ($3\cdot2\cdot1=6$). Άρτιοι είναι οι $132$ και $312$. Πιθανότητα $\dfrac{2}{6}=\dfrac{1}{3}$.
\solr{Γ3}{15} α) Σωστό. β) Λάθος: ισοπίθανα είναι τα $36$ ζευγάρια. γ) Σωστό. δ) Σωστό. ε) Λάθος: $2^4=16$.
\solr{Γ3}{16} α) Σωστό: $3\cdot2=6$. β) Λάθος: πλησιάζει την πιθανότητα όσο αυξάνονται οι επαναλήψεις. γ) Σωστό. δ) Σωστό. ε) Λάθος: δεν είναι σίγουρο ότι θα συμβεί ακριβώς $50$ φορές.
